Let $p_1=2<p_2=3<p_3=5<\cdots$ be the increasing sequence of primes and define the prime gaps
\[
  d_n := p_{n+1}-p_n \qquad (n\ge 1).
\]
\emph{Ambiguity note.} The problem text says: ``$h(x)$ be maximal such that for some $n<x$ the numbers $d_n,\ldots,d_{n+h(x)-1}$ are all distinct.''  As written, the endpoint index $n+h(x)-1$ is not constrained to be $\le x$.

\emph{Minimal corrected convention used below.} For an integer $x\ge 1$, define
\[
  h(x) := \max\Bigl\{h\in\mathbb{N}: \exists n\in\{1,2,\dots,x-h+1\}\ \text{s.t.}\ d_n,d_{n+1},\dots,d_{n+h-1}\ \text{are pairwise distinct}\Bigr\}.
\]
Equivalently, $h(x)$ is the maximum length of a contiguous block of pairwise distinct values inside the finite sequence $(d_1,\dots,d_x)$.

The task is to estimate $h(x)$ as $x\to\infty$, and in particular to decide whether:
\[
  (\exists c>0)(\forall x\ \text{large})\quad h(x)>(\log x)^c,\qquad\text{and whether }\quad h(x)=o(\log x).
\]
